% !TEX root = ../../main.tex
% C6.1 — Đánh giá và lựa chọn công nghệ, mô hình (RÚT GỌN 2026-09-29; bản cũ ở report/archive/)
% Phiên bản lấy từ frontend/package.json, backend/pyproject.toml + requirements lock, infra/compose.yaml.
% Lịch sử lựa chọn: architect.md mục 9, 12 (FastAPI + pgvector + tệp cục bộ được gợi ý ban đầu -> Flask + Milvus + MinIO + PostgreSQL;
% RaSa theo đề xuất của giảng viên; CLIP đa ngôn ngữ từng được cân nhắc; OpenVINO để sau). AI plan: YOLOX dự phòng giấy phép (đăng ký, chưa có trọng số).
% Số liệu YOLO11 lấy từ docs.ultralytics.com/models/yolo11 (đã kiểm chứng). pgvector: lọc sau khi quét chỉ mục (README chính thức).
% 3.2 dẫn sec:6.1.4 (lý do chọn YOLO11n); 3.6 dẫn sec:6.1.3 (lý do chọn Milvus).
\section{Đánh giá và lựa chọn công nghệ, mô hình}
\label{sec:6.1}

Các lựa chọn được đánh giá theo bốn tiêu chí: yêu cầu ở Chương~\ref{chapter:phantich}; chạy được trên máy chỉ có CPU và 16~GB bộ nhớ; mã nguồn mở với giấy phép phù hợp đồ án phi thương mại; và mức quen thuộc của tác giả, để giảm rủi ro. Bảng~\ref{tab:tech-summary} tóm tắt các lựa chọn; các lựa chọn chính được giải thích dưới đây.

\begin{table}[H]
\centering
\caption{Tóm tắt các lựa chọn công nghệ và mô hình}
\label{tab:tech-summary}
\begin{tabularx}{\textwidth}{|>{\raggedright\arraybackslash}p{3.0cm}|>{\raggedright\arraybackslash}p{3.1cm}|>{\raggedright\arraybackslash}p{3.4cm}|L|}
\hline
\textbf{Thành phần} & \textbf{Lựa chọn} & \textbf{Phương án đã cân nhắc} & \textbf{Lý do chính} \\
\hline
Giao diện & React 19, Vite, Tailwind CSS & Sinh trang ở máy chủ; Angular, Vue & Nhiều tương tác trên một màn hình; tách giao diện và máy chủ qua API \\
\hline
Máy chủ & Python, Flask 3.1, SQLAlchemy, Alembic & Django, FastAPI & Cùng ngôn ngữ với các mô hình AI; gọn, tường minh \\
\hline
Hàng đợi & Bảng trong PostgreSQL & Celery kèm Redis hoặc RabbitMQ & Một tiến trình xử lý tuần tự; ít thành phần hạ tầng \\
\hline
CSDL quan hệ & PostgreSQL 17 & MySQL & Ràng buộc CHECK, chỉ mục có điều kiện, JSONB, \texttt{SKIP LOCKED} \\
\hline
Lưu trữ vector & Milvus 2.6 & pgvector, FAISS & Lọc trước khi tìm, lưu trữ bền vững, tách collection theo phiên bản \\
\hline
Lưu trữ khung hình & MinIO & Tệp trên ổ đĩa & Phân quyền theo bucket, mã băm, dùng chung với Milvus \\
\hline
Phát hiện người & YOLO11n & YOLO11s, YOLO11m, YOLOX, Faster R-CNN & Nhanh nhất trên CPU; YOLOX dự phòng về giấy phép \\
\hline
Theo vết & ByteTrack (mặc định), BoT-SORT & DeepSORT, StrongSORT & Không cần đặc trưng ngoại hình cho mỗi người mỗi khung hình \\
\hline
Bộ mã hóa & RaSa & Hai mô hình riêng; CLIP & Một không gian chung cho ba hình thức truy vấn \\
\hline
\end{tabularx}
\end{table}

\subsection{Công nghệ phía giao diện}
\label{sec:6.1.1}

Giao diện có nhiều tương tác trên một màn hình, như kéo thả ảnh truy vấn, sắp xếp lại kết quả hay mở hộp thoại chồng lên danh sách, nên phù hợp hướng \emph{ứng dụng một trang}: tải một lần rồi chỉ trao đổi dữ liệu qua API. Đề tài chọn React nhờ hệ sinh thái thư viện điều hướng, gọi API và định kiểu, cùng Vite để đóng gói và Tailwind CSS cho bố cục thống nhất, co giãn trên điện thoại.

\subsection{Công nghệ phía máy chủ}
\label{sec:6.1.2}

Máy chủ và tiến trình nền viết bằng Python, ngôn ngữ của mọi mô hình trong đề tài, nên máy chủ nạp trực tiếp được bộ mã hóa truy vấn và dùng chung mã truy cập dữ liệu với tiến trình nền. Django nặng hơn nhu cầu của một máy chủ chỉ có API, còn lợi thế bất đồng bộ của FastAPI không đáng kể khi phân tích video đã tách khỏi máy chủ (Mục~\ref{sec:5.1.3}); Flask gọn, tường minh và tác giả đã quen dùng. Với một tiến trình xử lý tuần tự, một bảng PostgreSQL thay cho Celery cùng Redis hay RabbitMQ, tiết kiệm bộ nhớ trên máy triển khai.

\subsection{Hệ quản trị cơ sở dữ liệu và lưu trữ}
\label{sec:6.1.3}

PostgreSQL có những tính năng thiết kế dùng trực tiếp: ràng buộc CHECK phức tạp, chỉ mục có điều kiện, JSONB cho nhật ký và số đo, và \texttt{SKIP LOCKED} để nhận công việc an toàn. Với vector, pgvector~\cite{pgvector2026} chỉ cần một cơ sở dữ liệu nhưng với chỉ mục gần đúng thì lọc \emph{sau} khi quét nên có thể trả thiếu kết quả (Mục~\ref{sec:3.6.3}); FAISS là thư viện trong bộ nhớ, để ứng dụng tự lo lưu trữ, lọc và truy cập đồng thời. Milvus~\cite{wang2021milvus} lọc trước khi tìm~\cite{milvus2026filtered}, đáp ứng trực tiếp phân quyền theo khu vực; nó cần thêm etcd và kho đối tượng, nhưng toàn bộ dịch vụ trong Docker chỉ dùng khoảng 515~MiB bộ nhớ (Mục~\ref{sec:8.7}). Khung hình lưu trong MinIO thay vì tệp thường để kiểm soát quyền theo bucket và lưu mã băm từng đối tượng, dùng chung máy chủ MinIO mà Milvus vốn cần, ở bucket riêng.

\subsection{Các mô hình AI}
\label{sec:6.1.4}

\paragraph{Phát hiện người.} Hướng hai giai đoạn như Faster R-CNN quá chậm trên CPU, nên đề tài chọn họ YOLO. Trong YOLO11, kích cỡ quyết định tốc độ trên CPU (Bảng~\ref{tab:yolo11-sizes}), và ngay YOLO11n đã khiến quy trình chậm hơn thời gian thực khoảng bảy lần (Mục~\ref{sec:7.2}), nên đề tài chọn kích cỡ nhỏ nhất, bù độ chính xác bằng ngưỡng tin cậy thấp kết hợp theo vết. Giấy phép AGPL-3.0 của YOLO11 phù hợp đồ án phi thương mại nhưng không phù hợp sản phẩm, nên YOLOX~\cite{ge2021yolox} (Apache-2.0) được đăng ký làm dự phòng.

\begin{table}[H]
\centering
\caption{So sánh các kích cỡ YOLO11 trên tập kiểm định COCO~\cite{ultralytics2024yolo11docs}}
\label{tab:yolo11-sizes}
\begin{tabularx}{\textwidth}{|L|>{\raggedright\arraybackslash}p{2.6cm}|>{\raggedright\arraybackslash}p{3.8cm}|>{\raggedright\arraybackslash}p{3.0cm}|}
\hline
\textbf{Mô hình} & \textbf{mAP$_{50\text{--}95}$} & \textbf{Thời gian trên CPU (ms/ảnh)} & \textbf{Số tham số (triệu)} \\
\hline
YOLO11n & 39,5 & 56,1 & 2,6 \\
\hline
YOLO11s & 47,0 & 90,0 & 9,4 \\
\hline
YOLO11m & 51,5 & 183,2 & 20,1 \\
\hline
\end{tabularx}
\end{table}

\paragraph{Theo vết.} Các thuật toán kết hợp đặc trưng ngoại hình như DeepSORT~\cite{wojke2017deepsort} và StrongSORT~\cite{du2023strongsort} giữ track tốt hơn khi bị che khuất lâu nhưng phải chạy thêm một mô hình cho mỗi người trên mỗi khung hình, quá nặng cho CPU ở cảnh đông người. Với camera cố định, ByteTrack~\cite{zhang2022bytetrack} là mặc định, còn BoT-SORT~\cite{aharon2022botsort} là phương án thay thế với ReID và bù chuyển động camera được tắt.

\paragraph{Bộ mã hóa.} Hai mô hình riêng cho ảnh và văn bản sẽ cần hai vector mỗi lần xuất hiện và hai mô hình khi lập chỉ mục, còn mô hình tổng quát như CLIP~\cite{radford2021clip}, kể cả bản đa ngôn ngữ cho tiếng Việt, kém phân biệt chi tiết nhỏ về trang phục. RaSa~\cite{bai2023rasa}, theo đề xuất của giảng viên hướng dẫn, phục vụ cả ba hình thức truy vấn với một vector mỗi lần xuất hiện (Mục~\ref{sec:3.5.4}), đổi lại chỉ hỗ trợ tiếng Anh; Mục~\ref{sec:8.4} đánh giá nó trên dữ liệu của đề tài.
